\documentclass[11pt, letterpaper]{article}

% --- Idioma y codificación ---
\usepackage[spanish]{babel}
\usepackage[utf8]{inputenc}
\usepackage[T1]{fontenc}
\usepackage{lmodern}

% --- Matemáticas ---
\usepackage{amsmath, amssymb, amsthm}

% --- Formato ---
\usepackage[margin=2.5cm]{geometry}
\usepackage{parskip}
\usepackage{booktabs}
\usepackage{hyperref}
\hypersetup{colorlinks=true, linkcolor=blue!70!black, urlcolor=blue!70!black}

% --- Entornos de teoremas ---
\theoremstyle{definition}
\newtheorem{definition}{Definición}[section]
\theoremstyle{plain}
\newtheorem{theorem}{Teorema}[section]
\newtheorem{lemma}[theorem]{Lema}
\newtheorem{corollary}[theorem]{Corolario}
\theoremstyle{remark}
\newtheorem{remark}{Observación}[section]

\title{%
  T01 --- Demostraciones formales:\\[4pt]
  Recursión, inducción y Fibonacci
}
\author{Curso Linux--C--Rust --- B15 Estructuras de datos en C}
\date{}

\begin{document}
\maketitle
\tableofcontents

% ===================================================================
\section{Recordatorio: inducción matemática}
% ===================================================================

\begin{definition}[Inducción simple]
Para probar $P(n)$ para todo $n \geq 0$:
\begin{enumerate}
  \item \textbf{Base}: probar $P(0)$.
  \item \textbf{Paso}: suponer $P(k)$ (hipótesis inductiva), probar $P(k+1)$.
\end{enumerate}
\end{definition}

\begin{definition}[Inducción fuerte (o completa)]
Para probar $P(n)$ para todo $n \geq 0$:
\begin{enumerate}
  \item \textbf{Base}: probar $P(0)$ (y posiblemente $P(1), \ldots, P(j)$).
  \item \textbf{Paso}: suponer $P(0), P(1), \ldots, P(k)$ \emph{todas}
    verdaderas, probar $P(k+1)$.
\end{enumerate}
\end{definition}

La inducción fuerte es más poderosa cuando el caso recursivo depende
de más de un predecesor (como Fibonacci, que depende de $F(n-1)$ y
$F(n-2)$).

% ===================================================================
\section{Demostración 1: correctitud del factorial recursivo}
% ===================================================================

\begin{definition}[Factorial recursivo]
\[
  f(n) = \begin{cases} 1 & \text{si } n = 0 \\ n \cdot f(n-1) & \text{si } n \geq 1 \end{cases}
\]
\end{definition}

\begin{definition}[Factorial por productoria]
\[
  n! = \prod_{i=1}^{n} i = 1 \cdot 2 \cdot 3 \cdots n
\]
con la convención de que el producto vacío ($n = 0$) es~1.
\end{definition}

\begin{theorem}\label{thm:factorial}
$f(n) = n!$ para todo $n \geq 0$.
\end{theorem}

\begin{proof}
Por inducción fuerte. Sea $P(n)$: ``$f(n) = n!$''.

\textbf{Base ($n = 0$).}
\[
  f(0) = 1 = 0! \quad \checkmark
\]

\textbf{Paso inductivo.} Sea $k \geq 0$. Suponemos $P(i)$ verdadera
para todo $0 \leq i \leq k$. Calculamos $f(k+1)$:
\begin{align*}
  f(k+1) &= (k+1) \cdot f(k)       && \text{(definición recursiva)} \\
         &= (k+1) \cdot k!          && \text{(hipótesis inductiva)} \\
         &= (k+1)!                   && \text{(definición de factorial)}
\end{align*}

$P(k+1)$ es verdadera. Por inducción fuerte, $f(n) = n!$ para todo
$n \geq 0$.
\end{proof}

\begin{remark}
Esta demostración también funciona con inducción simple, ya que el
caso recursivo solo depende de $f(n-1)$. Usamos inducción fuerte por
consistencia con la Demostración~3.
\end{remark}

% ===================================================================
\section{Demostración 2: terminación del factorial recursivo}
% ===================================================================

\begin{theorem}\label{thm:termination}
Para todo $n \geq 0$, la evaluación de $f(n)$ termina en exactamente
$n$ llamadas recursivas.
\end{theorem}

\begin{proof}
Por inducción simple.

\textbf{Base ($n = 0$).} $f(0) = 1$ retorna directamente. $0$
llamadas recursivas. \checkmark

\textbf{Paso inductivo.} Supongamos que $f(k)$ termina en $k$
llamadas. Entonces $f(k+1) = (k+1) \cdot f(k)$ ejecuta una llamada
a $f(k)$, que termina en $k$ llamadas. Total: $1 + k = k+1$ llamadas.
\end{proof}

\begin{remark}[Función de decrecimiento]
El argumento $n$ decrece estrictamente en cada llamada ($n \to n-1$)
y está acotado inferiormente por~0. Todo entero no negativo
decreciente eventualmente llega a~0, donde la recursión se detiene.
\end{remark}

% ===================================================================
\section{Demostración 3: fórmula cerrada de Fibonacci}
% ===================================================================

\begin{definition}[Secuencia de Fibonacci]
\[
  F(n) = \begin{cases}
    0 & \text{si } n = 0 \\
    1 & \text{si } n = 1 \\
    F(n-1) + F(n-2) & \text{si } n \geq 2
  \end{cases}
\]
\end{definition}

\begin{theorem}[Fórmula de Binet]\label{thm:binet}
Para todo $n \geq 0$:
\[
  F(n) = \frac{\varphi^n - \psi^n}{\sqrt{5}}
\]
donde $\varphi = \dfrac{1 + \sqrt{5}}{2}$ (proporción áurea) y
$\psi = \dfrac{1 - \sqrt{5}}{2}$.
\end{theorem}

\subsection{Parte 1: derivación por ecuación característica}

La recurrencia $F(n) = F(n-1) + F(n-2)$ es lineal homogénea de
segundo orden. Buscamos soluciones de la forma $F(n) = r^n$:
\[
  r^n = r^{n-1} + r^{n-2}
  \quad\Longrightarrow\quad
  r^2 - r - 1 = 0
\]

Resolviendo:
\[
  r = \frac{1 \pm \sqrt{5}}{2}
  \quad\Longrightarrow\quad
  \varphi = \frac{1 + \sqrt{5}}{2} \approx 1.618,
  \quad
  \psi = \frac{1 - \sqrt{5}}{2} \approx -0.618
\]

La solución general es $F(n) = A\varphi^n + B\psi^n$. Las condiciones
iniciales determinan $A$ y $B$:
\begin{align*}
  F(0) = 0: \quad & A + B = 0 \tag{I} \\
  F(1) = 1: \quad & A\varphi + B\psi = 1 \tag{II}
\end{align*}

De (I): $B = -A$. Sustituyendo en (II):
\[
  A(\varphi - \psi) = 1
  \quad\Longrightarrow\quad
  A \cdot \sqrt{5} = 1
  \quad\Longrightarrow\quad
  A = \frac{1}{\sqrt{5}}, \quad B = -\frac{1}{\sqrt{5}}
\]
donde usamos $\varphi - \psi = \sqrt{5}$.

\subsection{Parte 2: verificación por inducción fuerte}

\begin{proof}[Prueba del Teorema~\ref{thm:binet}]
Sea $P(n)$: ``$F(n) = (\varphi^n - \psi^n)/\sqrt{5}$''.

\textbf{Base ($n = 0$ y $n = 1$).}
\begin{align*}
  F(0) &= \frac{\varphi^0 - \psi^0}{\sqrt{5}} = \frac{1 - 1}{\sqrt{5}} = 0 \quad \checkmark \\
  F(1) &= \frac{\varphi - \psi}{\sqrt{5}} = \frac{\sqrt{5}}{\sqrt{5}} = 1 \quad \checkmark
\end{align*}

\textbf{Paso inductivo.} Sea $k \geq 1$. Suponemos $P(i)$ verdadera
para todo $0 \leq i \leq k$. Por la recurrencia y la hipótesis:
\[
  F(k+1) = F(k) + F(k-1)
  = \frac{\varphi^k - \psi^k}{\sqrt{5}}
  + \frac{\varphi^{k-1} - \psi^{k-1}}{\sqrt{5}}
  = \frac{\varphi^{k-1}(\varphi + 1) - \psi^{k-1}(\psi + 1)}{\sqrt{5}}
\]

Como $\varphi$ y $\psi$ son raíces de $r^2 = r + 1$:
\[
  \varphi + 1 = \varphi^2, \qquad \psi + 1 = \psi^2
\]

Verificación: $\varphi^2 = \frac{(1+\sqrt{5})^2}{4} = \frac{6+2\sqrt{5}}{4} = \frac{3+\sqrt{5}}{2} = \varphi + 1$. \checkmark

Sustituyendo:
\[
  F(k+1) = \frac{\varphi^{k-1} \cdot \varphi^2 - \psi^{k-1} \cdot \psi^2}{\sqrt{5}}
  = \frac{\varphi^{k+1} - \psi^{k+1}}{\sqrt{5}}
\]

$P(k+1)$ es verdadera. Por inducción fuerte, la fórmula vale para
todo $n \geq 0$.
\end{proof}

% ===================================================================
\section{Consecuencias de la fórmula de Binet}
% ===================================================================

\begin{corollary}[$F(n)$ siempre es entero]\label{cor:integer}
A pesar de que la fórmula involucra $\sqrt{5}$, $F(n) \in \mathbb{Z}$
para todo $n \geq 0$.
\end{corollary}

\begin{proof}
Expandiendo con el binomio de Newton:
\[
  \varphi^n = \frac{1}{2^n}\sum_{j=0}^{n}\binom{n}{j}(\sqrt{5})^j,
  \qquad
  \psi^n = \frac{1}{2^n}\sum_{j=0}^{n}\binom{n}{j}(-\sqrt{5})^j
\]

Al restar, los términos con $j$ par se cancelan y los impares se
duplican. Cada término con $j$ impar contiene
$(\sqrt{5})^j = 5^{(j-1)/2}\sqrt{5}$, que al dividir por $\sqrt{5}$
da un racional. Los coeficientes binomiales garantizan que el
resultado es entero.
\end{proof}

\begin{corollary}\label{cor:round}
$F(n) = \operatorname{round}(\varphi^n / \sqrt{5})$ para todo $n \geq 1$.
\end{corollary}

\begin{proof}
Como $|\psi| = \frac{\sqrt{5}-1}{2} \approx 0.618 < 1$:
\[
  \left|F(n) - \frac{\varphi^n}{\sqrt{5}}\right|
  = \frac{|\psi|^n}{\sqrt{5}}
  < \frac{1}{\sqrt{5}}
  < \frac{1}{2}
\]

$F(n)$ es entero y $\varphi^n/\sqrt{5}$ está a distancia $< 1/2$,
así que $F(n)$ es el entero más cercano.
\end{proof}

\begin{corollary}[Complejidad de la recursión ingenua]\label{cor:complexity}
El número de llamadas $T(n)$ de la implementación recursiva directa
satisface $T(n) = \Theta(\varphi^n)$.
\end{corollary}

\begin{proof}
$T(n)$ satisface $T(0) = T(1) = 1$, $T(n) = T(n-1) + T(n-2) + 1$.
Se puede probar que $T(n) = 2F(n+1) - 1$, lo que por la fórmula de
Binet da $T(n) = \Theta(\varphi^n)$.
\end{proof}

% ===================================================================
\section{Resumen de resultados}
% ===================================================================

\begin{table}[h]
\centering
\begin{tabular}{llc}
\toprule
\textbf{Resultado} & \textbf{Técnica} & \textbf{Tipo de inducción} \\
\midrule
$f(n) = n!$
  & Inducción
  & Simple (basta $P(k)$) \\
$f(n)$ termina en $n$ llamadas
  & Inducción + func.\ de decrecimiento
  & Simple \\
Fórmula de Binet para $F(n)$
  & Ec.\ característica + inducción
  & Fuerte ($P(k)$ y $P(k{-}1)$) \\
$F(n) \in \mathbb{Z}$
  & Expansión binomial
  & Directa \\
$F(n) = \operatorname{round}(\varphi^n/\sqrt{5})$
  & Cota sobre $|\psi^n|$
  & Directa \\
\bottomrule
\end{tabular}
\end{table}

\end{document}
